\begingroup
\setlength{\tabcolsep}{4.5pt}
\renewcommand{\arraystretch}{1.28}
\begin{longtable}{@{}>{\raggedright\arraybackslash}p{.15\textwidth}>{\raggedright\arraybackslash}p{.31\textwidth}>{\raggedright\arraybackslash}p{.25\textwidth}>{\raggedright\arraybackslash}p{.23\textwidth}@{}}
\caption{Sectores topológicos y virtuales: invariantes nativos de fusión, trenzado y persistencia.}\label{tab:masas-invariantes-topologicos-virtuales-rev3}\\
\toprule
\textbf{Sector} & \textbf{Objeto y relaciones} & \textbf{Invariante nativo} & \textbf{Alcance físico} \\
\midrule
\endfirsthead
\toprule
\textbf{Sector} & \textbf{Objeto y relaciones} & \textbf{Invariante nativo} & \textbf{Alcance físico} \\
\midrule
\endhead
\midrule
\multicolumn{4}{r}{\footnotesize Continúa en la página siguiente.}\\
\endfoot
\bottomrule
\endlastfoot
Fibonacci & $K=\mathbb Z1\oplus\mathbb Z\tau$, $\tau^2=1+\tau$. & $\operatorname{FPdim}(\tau)=\varphi$ en el anillo basado. & Las matrices $F/R$, pentágono, hexágono, Hamiltoniano y brecha pertenecen a la realización posterior. \\
Ising & $K=\mathbb Z\{1,\psi,\sigma\}$, $\psi^2=1$, $\psi\sigma=\sigma$, $\sigma^2=1+\psi$. & $\operatorname{FPdim}(\sigma)=\sqrt{2}$; representación de trenzado. & La realización Majorana conserva su propio morfismo físico. \\
$\mathbb Z/6$ & $q_{\epsilon,t}=(-1)^\epsilon\omega_3^t=\zeta_6^k$. & Carácter cuadrático y representación de trenzado sobre los seis sectores. & La observación pertinente es topológica, no una masa universal. \\
Secciones virtuales & Sistema inverso de prolongaciones con horizonte $L(s_N)$. & Registro finito, firma, acarreo y clase de obstrucción. & El observable nativo es la persistencia de la prolongación. \\
\end{longtable}
\endgroup
